\clearpage
\section{계산 절차와 구조 요약}
\label{sec:equation-group-workflow}

\begin{learninggoal}
구체적인 다항식의 갈루아군을 계산하는 표준 절차를 정리한다. 이차·삼차·특별한 사차식에서 인수분해, 판별식, 분해체 차수와 순열부분군을 연결한다.
\end{learninggoal}

\subsection{갈루아군 계산의 기본 절차}

분리다항식 $f\in K[X]$가 주어졌다고 하자. 다음 순서가 계산의 기본 골격이다.

\begin{enumerate}
  \item \textbf{분리성을 확인한다.} $\gcd(f,f')=1$인지 검사한다.
  \item \textbf{$K$ 위 인수분해를 구한다.} 기약인수의 차수는 근 집합의 갈루아 궤도 크기이다.
  \item \textbf{분해체를 묘사한다.} 필요한 근이나 단위근을 차례로 첨가한다.
  \item \textbf{분해체 차수를 구한다.} 이 수가 갈루아군의 위수이다.
  \item \textbf{$S_n$의 가능한 부분군을 좁힌다.} 기약이면 추이성, 판별식이 제곱이면 $A_n$ 포함 조건 등을 사용한다.
  \item \textbf{실제 자동동형을 구성한다.} 생성원의 상과 관계식을 적고 군의 표시를 확인한다.
\end{enumerate}

\begin{keyidea}
계산은 ``군을 추측한 뒤 이름을 붙이는 일''이 아니다. 먼저 분해체 차수로 군의 위수를 고정하고, 근에 대한 작용과 생성원 관계로 군구조를 확정해야 한다.
\end{keyidea}

\subsection{삼차다항식의 완전한 빠른 판정}

$\charac K\ne2,3$이고 $f\in K[X]$가 분리 삼차다항식이라 하자.

\begin{enumerate}
  \item $f$가 세 일차인수로 분해되면 갈루아군은 자명하다.
  \item $f$가 일차인수와 기약 이차인수의 곱이면 분해체는 이차확장이고 갈루아군은 $C_2$이다.
  \item $f$가 기약이면 갈루아군은 $C_3$ 또는 $S_3$이다.
  \item 기약인 경우 판별식이 $K$의 제곱이면 $C_3$, 아니면 $S_3$이다.
\end{enumerate}

\begin{example}
\[
  f=X^3-4X=X(X-2)(X+2)
\]
는 $\Q$에서 완전히 분해되므로 갈루아군은 자명하다.
\end{example}

\begin{example}
\[
  f=(X-1)(X^2-2)
\]
의 분해체는 $\Q(\sqrt2)$이고
\[
  \Gal(f/\Q)\cong C_2.
\]
근 집합의 궤도 크기는 $1$과 $2$이다.
\end{example}

\begin{example}
\[
  f=X^3-3X+1
\]
은 기약이고 판별식이 $81$이므로 갈루아군은 $C_3$이다.
\end{example}

\begin{example}
\[
  f=X^3-X+1
\]
은 기약이고 판별식이 $-23$이므로 갈루아군은 $S_3$이다.
\end{example}

\subsection{군론적 정보가 주는 체의 정보}

다항식의 갈루아군 $G$를 알면 유한 갈루아 대응으로 분해체의 모든 중간체를 분류할 수 있다.

\begin{itemize}
  \item $G=C_p$이면 진중간체가 없다.
  \item $G=V_4$이면 세 개의 이차 중간체가 있다.
  \item $G=S_3$이면 세 개의 비정규 삼차 중간체와 하나의 정규 이차 중간체가 있다.
  \item $G=D_4$이면 여러 이차·사차 중간체가 나타나며, 정규부분군에 대응하는 중간체만 기준체 위에서 갈루아이다.
\end{itemize}

특히 기약 삼차식의 $S_3$ 경우에는 판별식의 제곱근이 유일한 이차 중간체를 생성한다.

\subsection{이 장의 한계와 다음 장}

이 장의 방법만으로 모든 사차·오차다항식의 갈루아군을 효율적으로 결정할 수 있는 것은 아니다. 다음 도구가 더 필요하다.

\begin{itemize}
  \item 판별식의 일반 정의와 계수 공식,
  \item 두 다항식의 공통근을 판정하는 종결식,
  \item 사차다항식의 삼차분해식,
  \item 소수에서의 인수분해형과 갈루아군의 순환형 사이의 관계,
  \item 원분체와 순환확장의 구체적 구조.
\end{itemize}

\begin{chapterreview}
\begin{itemize}
  \item 분리 $n$차다항식의 갈루아군은 $S_n$의 부분군이다.
  \item 분해체 차수는 갈루아군의 위수이고 $n!$을 나눈다.
  \item 기약인수의 근 집합은 정확히 갈루아 궤도이다.
  \item 다항식이 기약일 필요충분조건은 갈루아군이 모든 근 위에서 추이적으로 작용하는 것이다.
  \item 기약 소수차 다항식의 갈루아군은 그 소수 길이의 순환을 포함한다.
  \item 기약 삼차식의 갈루아군은 $C_3$ 또는 $S_3$이다.
  \item 삼차 판별식이 제곱이면 $C_3$, 제곱이 아니면 $S_3$이다.
  \item 특별한 기약 사차식에서는 분해체 차수와 근 관계에 따라 $D_4$ 또는 $V_4$ 등이 나타난다.
  \item 계수가 독립변수인 일반 $n$차방정식의 갈루아군은 $S_n$이다.
  \item 대칭 유리함수의 고정체는 기본대칭식으로 생성된다.
\end{itemize}
\end{chapterreview}

\subsection*{복습 카드}

\begin{itemize}
  \item \textbf{다항식의 갈루아군은?} 분리다항식의 분해체가 기준체 위에서 갖는 자동동형군.
  \item \textbf{기약성의 군론적 판정은?} 근 집합에 대한 작용의 추이성.
  \item \textbf{기약인수는 무엇에 대응하는가?} 근 집합의 갈루아 궤도.
  \item \textbf{기약 삼차식의 가능한 군은?} $C_3$ 또는 $S_3$.
  \item \textbf{삼차식에서 두 군을 구별하는 것은?} 판별식이 기준체의 제곱인지 여부.
  \item \textbf{일반 $n$차방정식의 군은?} 전체 대칭군 $S_n$.
  \item \textbf{기본대칭식의 역할은?} 근에 대칭인 정보를 계수로 표현한다.
\end{itemize}
